\chapter{Requirements Analysis}\label{chap:requirements}

\section{Introduction}

Requirements analysis establishes what a system must do and under which constraints before design begins
\cite{pressman2020se,sommerville2016software}. For \ProjectName{} it has a particular weight: the users are people
in distress and the staff who help them, so a missing requirement is not an inconvenience but a risk. This chapter
identifies the stakeholders, states the functional and non-functional requirements, expresses them as user
stories and use cases, records the business rules that keep the service safe and fair, and describes the overall
workflow. The requirements are stated once, here, and traced through the design and the tests by the matrix in
\cref{sec:traceability}.

\section{Stakeholders and System Actors}

Ten human roles exist in the system, plus the unauthenticated \role{Visitor} who can only read public pages
(\cref{tab:actors}). A user holds exactly one role at a time. Authentication is delegated to an identity provider,
but the role that decides what a person may do is stored in \ProjectName{}'s own database.

{\tabfont\rowband
\begin{longtable}{@{}>{\raggedright\arraybackslash}p{\dimexpr(\linewidth-4\tabcolsep)*178/1000\relax}>{\raggedright\arraybackslash}p{\dimexpr(\linewidth-4\tabcolsep)*822/1000\relax}@{}}
\caption{Human actors}\label{tab:actors}\\
\rowcolor{najdanavy}\textcolor{white}{\sffamily\bfseries\small Actor} & \textcolor{white}{\sffamily\bfseries\small Role in the platform}\\
\endfirsthead
\rowcolor{najdanavy}\textcolor{white}{\sffamily\bfseries\small Actor} & \textcolor{white}{\sffamily\bfseries\small Role in the platform}\\
\endhead
\bottomrule\endlastfoot
\textbf{Visitor} & Reads the public landing page, Terms, Privacy Policy, system status and support form; can start registration.\\
\textbf{Citizen} & Reports emergencies, tracks the response, chats with responders, keeps a medical profile, verifies their identity, applies to become a first responder.\\
\textbf{First Responder} & A vetted volunteer (an approved citizen). Keeps all citizen abilities and owns one personal response unit that they operate alone.\\
\textbf{Ambulance Crew} & Works shifts on an ambulance (LEAD or CREW); the LEAD runs missions and the hospital handoff. Linked to an ambulance station or a hospital.\\
\textbf{Police Officer} & Works shifts on a police car; linked to a police station.\\
\textbf{Firefighter} & Works shifts on a fire truck; linked to a fire station.\\
\textbf{Dispatcher} & Runs the live queue: reviews triage, resolves duplicates, assigns units, cancels, marks false reports, reviews caller history.\\
\textbf{Hospital Staff} & Sees incoming transfers and live patient vitals for their hospital, and the medical profile of patients inbound to it.\\
\textbf{Facility Manager} & Runs one facility: creates and retires its units and creates, disables and enables its staff. What they may create is decided by the facility type.\\
\textbf{Admin} & Everything a dispatcher can do, plus users, roles, facilities, units, applications, guidance and incident oversight. Cannot touch a Super Admin account.\\
\textbf{Super Admin} & Everything an Admin can do, plus legal-document publishing, identity-verification overrides, phone overrides, Firebase reconciliation, incident deletion and granting the Super Admin role.\\
\end{longtable}}


\subsection{How accounts are created}

Accounts arrive in four ways (\cref{tab:accounts}). Citizens register themselves; first responders are vetted
citizens whose application an administrator approves; and all other roles are provisioned by an administrator, or
by a facility manager for the staff of one facility. For provisioned accounts the server returns a one-time
password-reset link, so no password is ever shown, e-mailed or shared.

{\tabfont\rowband
\begin{longtable}{@{}>{\raggedright\arraybackslash}p{\dimexpr(\linewidth-6\tabcolsep)*189/1000\relax}>{\raggedright\arraybackslash}p{\dimexpr(\linewidth-6\tabcolsep)*568/1000\relax}>{\raggedright\arraybackslash}p{\dimexpr(\linewidth-6\tabcolsep)*243/1000\relax}@{}}
\caption{How accounts are created}\label{tab:accounts}\\
\rowcolor{najdanavy}\textcolor{white}{\sffamily\bfseries\small Role(s)} & \textcolor{white}{\sffamily\bfseries\small How the account is created} & \textcolor{white}{\sffamily\bfseries\small How the person signs in}\\
\endfirsthead
\rowcolor{najdanavy}\textcolor{white}{\sffamily\bfseries\small Role(s)} & \textcolor{white}{\sffamily\bfseries\small How the account is created} & \textcolor{white}{\sffamily\bfseries\small How the person signs in}\\
\endhead
\bottomrule\endlastfoot
\textbf{Citizen} & Self-registration. \textbf{E-mail path:} e-mail address verified by a six-digit code, then name, phone, address, gender, password and acceptance of the current Terms and Privacy Policy. \textbf{Google path:} sign in with Google, then complete phone, address and gender. & E-mail and password, or Google\\
\textbf{First Responder} & A citizen applies with documents; an Admin approves, which changes the role and creates the personal unit. & As a citizen\\
\textbf{Ambulance Crew, Police, Firefighter, Dispatcher, Hospital Staff, Admin, Super Admin} & Provisioned by an Admin (or by a Facility Manager for their own facility’s staff). The server creates the Firebase account with a throw-away password and returns a one-time password-reset link to hand to the employee. & E-mail and the password the employee chose through the link\\
\textbf{Facility Manager} & Provisioned by an Admin and linked to exactly one facility. & As an employee\\
\end{longtable}}


\subsection{External systems}

\ProjectName{} depends on a small number of external systems (\cref{tab:system-actors}). Each is used for one
purpose and is wrapped so that its failure degrades only its own feature (\cref{sec:failure}).

{\tabfont\rowband
\begin{longtable}{@{}>{\raggedright\arraybackslash}p{\dimexpr(\linewidth-4\tabcolsep)*216/1000\relax}>{\raggedright\arraybackslash}p{\dimexpr(\linewidth-4\tabcolsep)*784/1000\relax}@{}}
\caption{External systems}\label{tab:system-actors}\\
\rowcolor{najdanavy}\textcolor{white}{\sffamily\bfseries\small System actor} & \textcolor{white}{\sffamily\bfseries\small Role in the platform}\\
\endfirsthead
\rowcolor{najdanavy}\textcolor{white}{\sffamily\bfseries\small System actor} & \textcolor{white}{\sffamily\bfseries\small Role in the platform}\\
\endhead
\bottomrule\endlastfoot
\textbf{Firebase Authentication} & Identity provider: e-mail/\allowbreak{}password, Google and phone verification; issues the ID tokens the API verifies; carries the role as a custom claim.\\
\textbf{Firebase Cloud Messaging} & Delivers push notifications to Android and web clients.\\
\textbf{Supabase Storage} & Private object store for evidence media and responder-application documents, accessed only through short-lived signed URLs.\\
\textbf{OpenRouter (LLM)} & Language-model API used for report priority classification and duplicate-report detection.\\
\textbf{AI Model Service} & Internal FastAPI service that hosts versioned prediction models behind an API key (see \cref{sec:ai-components}).\\
\textbf{Didit} & Identity-verification provider; hosts the document-and-selfie check and reports the decision by signed webhook.\\
\textbf{Brevo} & Sends the registration verification codes by e-mail.\\
\textbf{Nominatim} & Reverse geocoding of report coordinates and backfilling of facility addresses.\\
\textbf{Overpass API} & Source of real facilities (hospitals, ambulance, police and fire stations) for seeding.\\
\textbf{OSRM and OpenFreeMap} & Driving routes and ETAs, and map tiles, used by the clients.\\
\textbf{EmailJS} & Delivers the public support form from the web console.\\
\textbf{GitHub Releases} & Hosts signed Android APKs; the API reports the latest release to the app.\\
\end{longtable}}


\subsection{Role capability matrix}

\Cref{tab:role-matrix} summarises what each role may do. Enforcement is on the server, by role and by ownership
(\cref{sec:authorisation}).

\begin{table}[!htbp]
\caption{Role capability matrix.}\label{tab:role-matrix}
\centering\tabfont\setlength{\tabcolsep}{4pt}\renewcommand{\arraystretch}{1.25}
\rowcolors{2}{najdapale}{white}
\begin{tabular}{@{}>{\raggedright\arraybackslash}p{3.7cm}*{8}{>{\centering\arraybackslash}p{1.12cm}}@{}}
\rowcolor{najdanavy}\textcolor{white}{\sffamily\bfseries\scriptsize Capability} &
 \textcolor{white}{\sffamily\bfseries\scriptsize Citizen} & \textcolor{white}{\sffamily\bfseries\scriptsize First resp.} &
 \textcolor{white}{\sffamily\bfseries\scriptsize Crew, police, fire} & \textcolor{white}{\sffamily\bfseries\scriptsize Dispatcher} &
 \textcolor{white}{\sffamily\bfseries\scriptsize Hospital staff} & \textcolor{white}{\sffamily\bfseries\scriptsize Facility mgr} &
 \textcolor{white}{\sffamily\bfseries\scriptsize Admin} & \textcolor{white}{\sffamily\bfseries\scriptsize Super admin} \\
Report an emergency, track it, chat & \yes & \yes & \scriptsize chat & \scriptsize chat & & & \scriptsize chat & \scriptsize chat \\
Own medical profile; apply to become a responder & \yes & \scriptsize profile & & & & & & \\
Run shifts, accept and work missions & & \yes & \yes & & & & & \\
Hospital handoff and vitals (lead of an ambulance) & & & \yes & & & & & \\
Incoming transfers and live vitals & & & & & \yes & & & \\
Live queue, assign, cancel, false report & & & & \yes & & & \yes & \yes \\
View a patient's medical profile (with consent) & & \scriptsize on mission & \scriptsize ambulance & \yes & \scriptsize inbound & & & \\
Manage own facility's units and staff & & & & & & \yes & & \\
Users, roles, facilities, units, applications, guidance & & & & & & & \yes & \yes \\
Super-admin functions & & & & & & & & \yes \\
\end{tabular}
\end{table}

\section{Functional Requirements}\label{sec:fr}

The functional requirements are grouped by capability. Each requirement has an identifier, a priority
(\emph{Must}, \emph{Should} or \emph{Could}) and the platform it applies to. All requirements listed are
implemented.

\subsection{Identity and access management}
{\tabfont\rowband
\begin{longtable}{@{}>{\raggedright\arraybackslash}p{\dimexpr(\linewidth-8\tabcolsep)*146/1000\relax}>{\raggedright\arraybackslash}p{\dimexpr(\linewidth-8\tabcolsep)*629/1000\relax}>{\raggedright\arraybackslash}p{\dimexpr(\linewidth-8\tabcolsep)*107/1000\relax}>{\raggedright\arraybackslash}p{\dimexpr(\linewidth-8\tabcolsep)*118/1000\relax}@{}}
\caption{Functional requirements: identity and access management}\label{tab:fr-iam}\\
\rowcolor{najdanavy}\textcolor{white}{\sffamily\bfseries\small ID} & \textcolor{white}{\sffamily\bfseries\small Requirement} & \textcolor{white}{\sffamily\bfseries\small Priority} & \textcolor{white}{\sffamily\bfseries\small Platform}\\
\endfirsthead
\rowcolor{najdanavy}\textcolor{white}{\sffamily\bfseries\small ID} & \textcolor{white}{\sffamily\bfseries\small Requirement} & \textcolor{white}{\sffamily\bfseries\small Priority} & \textcolor{white}{\sffamily\bfseries\small Platform}\\
\endhead
\bottomrule\endlastfoot
\textbf{FR-IAM-01} & A visitor can register as a citizen with an e-mail address. The system e-mails a six-digit verification code that expires after 10 minutes, can be re-sent only after a 60-second cool-down, and allows at most five wrong attempts. & Must & Web, Mobile\\
\textbf{FR-IAM-02} & Only after the e-mail is verified, the visitor supplies first and last name, phone number, address (governorate, city, street, optional landmark), gender and a password to create the account. Nothing is created in Firebase or the database before verification succeeds. & Must & Web, Mobile\\
\textbf{FR-IAM-03} & Registration is atomic across the identity provider and the database: if the database step fails, the newly created Firebase account is deleted again. & Must & Backend\\
\textbf{FR-IAM-04} & Code requests, code verification and registration are rate-limited per client address, and code requests are also limited per e-mail address. & Must & Backend\\
\textbf{FR-IAM-05} & A visitor can sign in with Google. If no NAJDA profile exists yet, the client collects phone, address and gender, and the account is created as a citizen. & Must & Web, Mobile\\
\textbf{FR-IAM-06} & The server alone decides whether a profile is complete (phone, governorate, city, street and gender present); clients use that flag to route the user to profile completion. & Must & Backend\\
\textbf{FR-IAM-07} & A signed-in user can sign in with e-mail and password, recover a forgotten password, and sign out. & Must & Web, Mobile\\
\textbf{FR-IAM-08} & A user can edit their own first name, last name, gender and address. & Must & Web, Mobile\\
\textbf{FR-IAM-09} & A user can change their e-mail address; the change is verified by Firebase and then synchronised into the NAJDA profile. & Should & Web, Mobile\\
\textbf{FR-IAM-10} & A user can verify a phone number with an SMS code or change it. A phone number already verified on another account cannot be claimed. & Must & Web, Mobile\\
\textbf{FR-IAM-11} & A user can change their password. & Must & Web, Mobile\\
\textbf{FR-IAM-12} & An Admin can provision accounts for staff roles. The server enforces the role--facility rules, creates the Firebase account with a temporary password and returns a one-time password-reset link; only a Super Admin can create another Super Admin. & Must & Web\\
\textbf{FR-IAM-13} & An Admin can search, filter (role, enabled) and page through all users, and edit a user’s name, gender and address. & Must & Web\\
\textbf{FR-IAM-14} & An Admin can change a user’s role. A plain Admin cannot modify a Super Admin or grant Super Admin; a Super Admin cannot remove their own Super Admin role. The new role is written to the Firebase custom claim, and a facility link that no longer fits the new role is released. & Must & Web\\
\textbf{FR-IAM-15} & An Admin can assign or clear a staff member’s facility, subject to the facility-type rules for the role. & Must & Web\\
\textbf{FR-IAM-16} & An Admin can disable and re-enable an account. The only remaining Super Admin cannot be disabled, and a plain Admin cannot disable a Super Admin. & Must & Web\\
\textbf{FR-IAM-17} & An Admin can generate a password-reset link for a user, mark a user’s phone as verified (audited), and override a user’s e-mail. A Super Admin can override a phone number without re-verification. & Should & Web\\
\textbf{FR-IAM-18} & A Super Admin can compare Firebase and the database, import a Firebase account that has no profile, remove an orphaned Firebase account, and repair a database user that has no Firebase account. & Should & Web\\
\textbf{FR-IAM-19} & A citizen with a verified e-mail and phone and a verified identity can apply to become a First Responder with a motivation and supporting documents; at most one application may be pending. & Must & Web, Mobile\\
\textbf{FR-IAM-20} & An Admin can list applications, open their documents through signed links, and approve or reject (with notes). Approval requires at least one document, changes the role to First Responder, updates the Firebase claim and creates the user’s personal response unit. & Must & Web\\
\end{longtable}}

\subsection{Identity verification, legal consent and medical profile}
{\tabfont\rowband
\begin{longtable}{@{}>{\raggedright\arraybackslash}p{\dimexpr(\linewidth-8\tabcolsep)*146/1000\relax}>{\raggedright\arraybackslash}p{\dimexpr(\linewidth-8\tabcolsep)*629/1000\relax}>{\raggedright\arraybackslash}p{\dimexpr(\linewidth-8\tabcolsep)*107/1000\relax}>{\raggedright\arraybackslash}p{\dimexpr(\linewidth-8\tabcolsep)*118/1000\relax}@{}}
\caption{Functional requirements: identity verification, legal consent and medical profile}\label{tab:fr-id}\\
\rowcolor{najdanavy}\textcolor{white}{\sffamily\bfseries\small ID} & \textcolor{white}{\sffamily\bfseries\small Requirement} & \textcolor{white}{\sffamily\bfseries\small Priority} & \textcolor{white}{\sffamily\bfseries\small Platform}\\
\endfirsthead
\rowcolor{najdanavy}\textcolor{white}{\sffamily\bfseries\small ID} & \textcolor{white}{\sffamily\bfseries\small Requirement} & \textcolor{white}{\sffamily\bfseries\small Priority} & \textcolor{white}{\sffamily\bfseries\small Platform}\\
\endhead
\bottomrule\endlastfoot
\textbf{FR-ID-01} & A user can start identity verification; the system opens a hosted Didit session and records the user’s verification status (not started, in progress, in review, verified, declined, expired, abandoned). & Must & Web, Mobile\\
\textbf{FR-ID-02} & The verification status is written only by the backend: from the signed Didit webhook, from a reconciliation with the provider when a status is stale, or from an audited Super Admin override. A person can start again after a decline, expiry or abandonment. & Must & Backend\\
\textbf{FR-ID-03} & Webhooks are authenticated by signature and timestamp, de-duplicated by event id, and confirmed against the provider’s decision before the status changes. & Must & Backend\\
\textbf{FR-ID-04} & Deployments can require verified identity before a citizen may report an emergency. The requirement can exempt roles and can grandfather accounts created before a cut-off date. & Should & Backend\\
\textbf{FR-ID-05} & A Super Admin can mark a user verified or reset a user’s verification (which also deletes the provider sessions). & Should & Web\\
\textbf{FR-LEG-01} & The system stores versioned Terms and Privacy Policy documents in English and Arabic. A Super Admin can publish a new version, optionally requiring re-acceptance. & Must & Web\\
\textbf{FR-LEG-02} & Registration records acceptance of the latest versions. When a version that requires acceptance is published, every user is notified and must accept before using protected features; the API answers 403 LEGAL\_\allowbreak{}CONSENT\_\allowbreak{}REQUIRED until they do. & Must & Web, Mobile\\
\textbf{FR-LEG-03} & Anyone can read the current Terms and Privacy Policy without signing in. & Must & Web, Mobile\\
\textbf{FR-MED-01} & A citizen or first responder can keep a medical profile: date of birth, weight, blood type, allergies, conditions (or an explicit “none”), medication, notes and up to three emergency contacts. & Should & Web, Mobile\\
\textbf{FR-MED-02} & Free-text medical fields and emergency-contact names and phones are encrypted at rest. & Must & Backend\\
\textbf{FR-MED-03} & Sharing the profile with responders is off by default and is a standing consent that can be withdrawn at any time; withdrawal takes effect immediately, even for open incidents. & Must & Web, Mobile\\
\textbf{FR-MED-04} & A report carries a per-report choice to share the profile. The server never attaches the reporter’s profile when the report is for someone else. & Must & Backend\\
\textbf{FR-MED-05} & Only people treating the patient can read the profile: a dispatcher, an ambulance or first-responder crew member on an active mission, and staff of the destination hospital; only while the incident is open. Administrators are excluded. Every read is audited and fails if the audit cannot be written. & Must & Backend\\
\end{longtable}}

\subsection{Emergency reporting}
{\tabfont\rowband
\begin{longtable}{@{}>{\raggedright\arraybackslash}p{\dimexpr(\linewidth-8\tabcolsep)*146/1000\relax}>{\raggedright\arraybackslash}p{\dimexpr(\linewidth-8\tabcolsep)*629/1000\relax}>{\raggedright\arraybackslash}p{\dimexpr(\linewidth-8\tabcolsep)*107/1000\relax}>{\raggedright\arraybackslash}p{\dimexpr(\linewidth-8\tabcolsep)*118/1000\relax}@{}}
\caption{Functional requirements: emergency reporting}\label{tab:fr-rep}\\
\rowcolor{najdanavy}\textcolor{white}{\sffamily\bfseries\small ID} & \textcolor{white}{\sffamily\bfseries\small Requirement} & \textcolor{white}{\sffamily\bfseries\small Priority} & \textcolor{white}{\sffamily\bfseries\small Platform}\\
\endfirsthead
\rowcolor{najdanavy}\textcolor{white}{\sffamily\bfseries\small ID} & \textcolor{white}{\sffamily\bfseries\small Requirement} & \textcolor{white}{\sffamily\bfseries\small Priority} & \textcolor{white}{\sffamily\bfseries\small Platform}\\
\endhead
\bottomrule\endlastfoot
\textbf{FR-REP-01} & A signed-in citizen can submit an emergency report with a category (medical, police or fire), a location, and an optional text message, in one action. & Must & Web, Mobile\\
\textbf{FR-REP-02} & The location comes from GPS or, where enabled by configuration, a manually placed pin that must lie within a configured radius (default 300 m) of the device’s own position; the location must be inside Egypt and is recorded with its source. & Must & Web, Mobile\\
\textbf{FR-REP-03} & Submission is idempotent: a client-generated request id guarantees that a retried request returns the original incident without creating a second one or repeating side effects. & Must & Backend\\
\textbf{FR-REP-04} & A citizen may have only one active incident at a time. & Must & Backend\\
\textbf{FR-REP-05} & A citizen can state that the emergency is for someone else; such a report never shares the reporter’s medical profile. & Should & Web, Mobile\\
\textbf{FR-REP-06} & A citizen can attach photos, video, audio and text as evidence, uploaded directly to private storage through signed URLs, and can add evidence while the incident is live. & Must & Web, Mobile\\
\textbf{FR-REP-07} & Until a unit is assigned, the citizen can edit the original text message and delete evidence. While the incident is active the citizen can update the number of injured people. & Should & Web, Mobile\\
\textbf{FR-REP-08} & A citizen can cancel their report until a unit has arrived, choosing a reason category. Open missions are stood down and the crews notified. & Must & Web, Mobile\\
\textbf{FR-REP-09} & A citizen can follow the incident live: status, assigned units with positions on a map, and an ETA. & Must & Web, Mobile\\
\textbf{FR-REP-10} & The citizen, dispatch staff and the crew of an active mission can exchange chat messages on the incident; chat content is encrypted at rest. & Must & Web, Mobile\\
\textbf{FR-REP-11} & A citizen can list and open their past incidents. & Must & Web, Mobile\\
\textbf{FR-REP-12} & The app presents bilingual first-aid and safety guidance cards for medical, police and fire situations, editable by administrators. & Should & Web, Mobile\\
\textbf{FR-REP-13} & The mobile app queues a report that cannot be sent (no connection) in a local outbox with its request id, and sends it automatically when connectivity returns. & Should & Mobile\\
\textbf{FR-REP-14} & The report address is resolved asynchronously by reverse geocoding and never delays the submission. & Should & Backend\\
\end{longtable}}

\subsection{Dispatch operations}
{\tabfont\rowband
\begin{longtable}{@{}>{\raggedright\arraybackslash}p{\dimexpr(\linewidth-8\tabcolsep)*146/1000\relax}>{\raggedright\arraybackslash}p{\dimexpr(\linewidth-8\tabcolsep)*629/1000\relax}>{\raggedright\arraybackslash}p{\dimexpr(\linewidth-8\tabcolsep)*107/1000\relax}>{\raggedright\arraybackslash}p{\dimexpr(\linewidth-8\tabcolsep)*118/1000\relax}@{}}
\caption{Functional requirements: dispatch operations}\label{tab:fr-dsp}\\
\rowcolor{najdanavy}\textcolor{white}{\sffamily\bfseries\small ID} & \textcolor{white}{\sffamily\bfseries\small Requirement} & \textcolor{white}{\sffamily\bfseries\small Priority} & \textcolor{white}{\sffamily\bfseries\small Platform}\\
\endfirsthead
\rowcolor{najdanavy}\textcolor{white}{\sffamily\bfseries\small ID} & \textcolor{white}{\sffamily\bfseries\small Requirement} & \textcolor{white}{\sffamily\bfseries\small Priority} & \textcolor{white}{\sffamily\bfseries\small Platform}\\
\endhead
\bottomrule\endlastfoot
\textbf{FR-DSP-01} & Dispatchers see a live queue of open incidents and a list of active incidents, updated in real time without refreshing. & Must & Web\\
\textbf{FR-DSP-02} & After submission the system classifies each report’s priority (critical, high, medium, low) with a confidence value using a language model; a failure never affects the incident, and a dispatcher can retry. & Must & Backend\\
\textbf{FR-DSP-03} & The system suggests a possible duplicate (same category, within 2 km and 120 minutes of an open incident). A dispatcher can confirm the duplicate or dismiss the suggestion. & Should & Web\\
\textbf{FR-DSP-04} & A dispatcher can see which colleagues are viewing an incident, to avoid double handling. & Could & Web\\
\textbf{FR-DSP-05} & A dispatcher can see live unit positions and statuses, ranked by distance and estimated travel time to the incident, and can offer a mission to an available unit. & Must & Web\\
\textbf{FR-DSP-06} & A unit can be offered a mission only if it is available, has an active LEAD on shift, and the incident is not closed. Concurrent assignments of one unit are serialised so a unit is never double-booked, and a unit is never sent to a report made by one of its own crew. & Must & Backend\\
\textbf{FR-DSP-07} & More than one unit can be assigned to the same incident. & Must & Backend\\
\textbf{FR-DSP-08} & A dispatcher can cancel a unit’s mission at any non-terminal stage; the unit becomes available. If it was the last active mission, the incident returns to dispatcher review. & Must & Web\\
\textbf{FR-DSP-09} & A dispatcher can complete an incident by override, which closes its remaining missions. & Should & Web\\
\textbf{FR-DSP-10} & A dispatcher can mark an incident as a false report (prank or hoax, spam or test, malicious). This closes the incident, stands units down and records the verdict against the caller; an Admin can clear the verdict. & Must & Web\\
\textbf{FR-DSP-11} & A dispatcher can review a caller’s history: previous reports, cancellations and false-report verdicts, and the resulting trust tier. & Should & Web\\
\textbf{FR-DSP-12} & A caller’s trust is assessed from recent history with decay (30-day half-life, 120-day window). The harshest automatic effect is a 10-minute wait between reports with the emergency numbers shown; reports are never refused outright. & Must & Backend\\
\textbf{FR-DSP-13} & An Admin sees a review list of callers whose trust tier warrants attention. & Should & Web\\
\textbf{FR-DSP-14} & A Super Admin can delete an incident. & Could & Web\\
\end{longtable}}

\subsection{Field response and shift management}
{\tabfont\rowband
\begin{longtable}{@{}>{\raggedright\arraybackslash}p{\dimexpr(\linewidth-8\tabcolsep)*146/1000\relax}>{\raggedright\arraybackslash}p{\dimexpr(\linewidth-8\tabcolsep)*629/1000\relax}>{\raggedright\arraybackslash}p{\dimexpr(\linewidth-8\tabcolsep)*107/1000\relax}>{\raggedright\arraybackslash}p{\dimexpr(\linewidth-8\tabcolsep)*118/1000\relax}@{}}
\caption{Functional requirements: field response and shift management}\label{tab:fr-fld}\\
\rowcolor{najdanavy}\textcolor{white}{\sffamily\bfseries\small ID} & \textcolor{white}{\sffamily\bfseries\small Requirement} & \textcolor{white}{\sffamily\bfseries\small Priority} & \textcolor{white}{\sffamily\bfseries\small Platform}\\
\endfirsthead
\rowcolor{najdanavy}\textcolor{white}{\sffamily\bfseries\small ID} & \textcolor{white}{\sffamily\bfseries\small Requirement} & \textcolor{white}{\sffamily\bfseries\small Priority} & \textcolor{white}{\sffamily\bfseries\small Platform}\\
\endhead
\bottomrule\endlastfoot
\textbf{FR-FLD-01} & An employee can start a shift on an available unit of their type and facility as LEAD, or join an already running shift as CREW; a first responder starts a shift on their personal unit. & Must & Web, Mobile\\
\textbf{FR-FLD-02} & A unit has at most one active LEAD. Starting a shift makes the unit available; when the LEAD ends the shift the unit goes offline, and a crew member leaving does not. & Must & Backend\\
\textbf{FR-FLD-03} & An employee can be on one shift at a time, and cannot end or leave a shift while the unit is on an active mission. & Must & Backend\\
\textbf{FR-FLD-04} & An Admin can force-end a unit’s shift. & Should & Web\\
\textbf{FR-FLD-05} & The unit’s crew receive an urgent push notification for each offered mission; the LEAD can accept it, or reject it with a reason. & Must & Web, Mobile\\
\textbf{FR-FLD-06} & The LEAD advances an accepted mission to en route and arrived; the citizen and dispatch see each step live. The LEAD can withdraw from an accepted or en-route mission. & Must & Web, Mobile\\
\textbf{FR-FLD-07} & The LEAD of a non-ambulance unit completes the mission on scene. Completing the last active mission resolves the incident. & Must & Web, Mobile\\
\textbf{FR-FLD-08} & A responder sees a route map to the incident and their teammates’ positions. & Should & Web, Mobile\\
\textbf{FR-FLD-09} & The LEAD’s device shares the unit’s live location; the LEAD can reset it to the home facility. Positions are visible to staff only. & Must & Web, Mobile\\
\textbf{FR-FLD-10} & A responder can review their mission history. & Should & Web, Mobile\\
\end{longtable}}

\subsection{Hospital coordination}
{\tabfont\rowband
\begin{longtable}{@{}>{\raggedright\arraybackslash}p{\dimexpr(\linewidth-8\tabcolsep)*146/1000\relax}>{\raggedright\arraybackslash}p{\dimexpr(\linewidth-8\tabcolsep)*629/1000\relax}>{\raggedright\arraybackslash}p{\dimexpr(\linewidth-8\tabcolsep)*107/1000\relax}>{\raggedright\arraybackslash}p{\dimexpr(\linewidth-8\tabcolsep)*118/1000\relax}@{}}
\caption{Functional requirements: hospital coordination}\label{tab:fr-hos}\\
\rowcolor{najdanavy}\textcolor{white}{\sffamily\bfseries\small ID} & \textcolor{white}{\sffamily\bfseries\small Requirement} & \textcolor{white}{\sffamily\bfseries\small Priority} & \textcolor{white}{\sffamily\bfseries\small Platform}\\
\endfirsthead
\rowcolor{najdanavy}\textcolor{white}{\sffamily\bfseries\small ID} & \textcolor{white}{\sffamily\bfseries\small Requirement} & \textcolor{white}{\sffamily\bfseries\small Priority} & \textcolor{white}{\sffamily\bfseries\small Platform}\\
\endhead
\bottomrule\endlastfoot
\textbf{FR-HOS-01} & Once an ambulance mission has arrived on scene, its LEAD can see registered hospitals within 30 km, nearest first, and select one as the destination. & Must & Web, Mobile\\
\textbf{FR-HOS-02} & Selecting a hospital notifies that hospital’s staff (urgent push and live update) that a patient is incoming. & Must & Backend\\
\textbf{FR-HOS-03} & The LEAD marks the transfer en route and arrived. Arrival completes the mission, frees the unit and resolves the incident when it is the last active mission. & Must & Web, Mobile\\
\textbf{FR-HOS-04} & While transporting, the LEAD can record vital signs (heart rate, SpO\textsubscript{2}, blood pressure, respiratory rate, temperature, consciousness level, notes); notes are encrypted at rest. Vitals are accepted only for a registered hospital destination. & Must & Web, Mobile\\
\textbf{FR-HOS-05} & Hospital staff see their hospital’s incoming transfers with a map, ETA and the live vitals stream and history. & Must & Web\\
\textbf{FR-HOS-06} & A destination that is not a registered facility can be recorded as free text, without vitals streaming. & Could & Web, Mobile\\
\end{longtable}}

\subsection{Facility, unit and content administration}
{\tabfont\rowband
\begin{longtable}{@{}>{\raggedright\arraybackslash}p{\dimexpr(\linewidth-8\tabcolsep)*146/1000\relax}>{\raggedright\arraybackslash}p{\dimexpr(\linewidth-8\tabcolsep)*629/1000\relax}>{\raggedright\arraybackslash}p{\dimexpr(\linewidth-8\tabcolsep)*107/1000\relax}>{\raggedright\arraybackslash}p{\dimexpr(\linewidth-8\tabcolsep)*118/1000\relax}@{}}
\caption{Functional requirements: facility, unit and content administration}\label{tab:fr-adm}\\
\rowcolor{najdanavy}\textcolor{white}{\sffamily\bfseries\small ID} & \textcolor{white}{\sffamily\bfseries\small Requirement} & \textcolor{white}{\sffamily\bfseries\small Priority} & \textcolor{white}{\sffamily\bfseries\small Platform}\\
\endfirsthead
\rowcolor{najdanavy}\textcolor{white}{\sffamily\bfseries\small ID} & \textcolor{white}{\sffamily\bfseries\small Requirement} & \textcolor{white}{\sffamily\bfseries\small Priority} & \textcolor{white}{\sffamily\bfseries\small Platform}\\
\endhead
\bottomrule\endlastfoot
\textbf{FR-ADM-01} & An Admin can create and edit facilities (hospital, ambulance station, fire station, police station) with name, address and coordinates, and list them with paging. & Must & Web\\
\textbf{FR-ADM-02} & Before changing a facility’s type the Admin sees an impact assessment (units and staff that would no longer fit). & Should & Web\\
\textbf{FR-ADM-03} & A first start loads a bundled list of Egyptian facilities (about 800). An Admin can additionally import hospitals, fire stations and police stations from OpenStreetMap through the Overpass API (ambulance stations are added manually) and backfill missing addresses through reverse geocoding; imports skip facilities already present. & Should & Web\\
\textbf{FR-ADM-04} & A facility is “registered” (usable as a destination or station) only while it has a linked active staff member or a non-archived unit; the flag is maintained automatically and reconciled nightly. & Must & Backend\\
\textbf{FR-ADM-05} & An Admin can create, edit and retire response units. A unit with history is archived instead of deleted, after a delete-impact check. & Must & Web\\
\textbf{FR-ADM-06} & An Admin can reconcile unit statuses that no longer match reality (for example a unit marked busy with no active mission). & Should & Web\\
\textbf{FR-ADM-07} & An Admin can search all incidents and edit the first-aid guidance content. & Should & Web\\
\textbf{FR-ADM-08} & A Facility Manager can list, create, edit and retire the units of their own facility, and list, create, disable and enable its staff. The unit type and the permitted staff roles are determined by the facility type, never by the request. & Must & Web\\
\end{longtable}}

\subsection{Notifications}
{\tabfont\rowband
\begin{longtable}{@{}>{\raggedright\arraybackslash}p{\dimexpr(\linewidth-8\tabcolsep)*146/1000\relax}>{\raggedright\arraybackslash}p{\dimexpr(\linewidth-8\tabcolsep)*629/1000\relax}>{\raggedright\arraybackslash}p{\dimexpr(\linewidth-8\tabcolsep)*107/1000\relax}>{\raggedright\arraybackslash}p{\dimexpr(\linewidth-8\tabcolsep)*118/1000\relax}@{}}
\caption{Functional requirements: notifications}\label{tab:fr-not}\\
\rowcolor{najdanavy}\textcolor{white}{\sffamily\bfseries\small ID} & \textcolor{white}{\sffamily\bfseries\small Requirement} & \textcolor{white}{\sffamily\bfseries\small Priority} & \textcolor{white}{\sffamily\bfseries\small Platform}\\
\endfirsthead
\rowcolor{najdanavy}\textcolor{white}{\sffamily\bfseries\small ID} & \textcolor{white}{\sffamily\bfseries\small Requirement} & \textcolor{white}{\sffamily\bfseries\small Priority} & \textcolor{white}{\sffamily\bfseries\small Platform}\\
\endhead
\bottomrule\endlastfoot
\textbf{FR-NOT-01} & Significant events create an in-app notification for each concerned user and a push notification to their registered devices; the inbox shows an unread count and supports mark-read and mark-all-read. & Must & Web, Mobile\\
\textbf{FR-NOT-02} & Notification text is produced in English and Arabic and delivered in the language of the recipient’s device. & Must & Backend\\
\textbf{FR-NOT-03} & Notifications are created only after the triggering transaction commits, so a rolled-back action never produces a notification. & Must & Backend\\
\textbf{FR-NOT-04} & Four notification types are urgent (new incident, mission offered, mission rejected, patient incoming) and use high-priority, short-lived push on a dedicated channel. & Should & Backend\\
\textbf{FR-NOT-05} & Devices register and unregister their push tokens; invalid tokens are pruned automatically. Notifications and tokens are purged on a schedule. & Must & Backend\\
\end{longtable}}

\subsection{Platform services}
{\tabfont\rowband
\begin{longtable}{@{}>{\raggedright\arraybackslash}p{\dimexpr(\linewidth-8\tabcolsep)*146/1000\relax}>{\raggedright\arraybackslash}p{\dimexpr(\linewidth-8\tabcolsep)*629/1000\relax}>{\raggedright\arraybackslash}p{\dimexpr(\linewidth-8\tabcolsep)*107/1000\relax}>{\raggedright\arraybackslash}p{\dimexpr(\linewidth-8\tabcolsep)*118/1000\relax}@{}}
\caption{Functional requirements: platform services}\label{tab:fr-plt}\\
\rowcolor{najdanavy}\textcolor{white}{\sffamily\bfseries\small ID} & \textcolor{white}{\sffamily\bfseries\small Requirement} & \textcolor{white}{\sffamily\bfseries\small Priority} & \textcolor{white}{\sffamily\bfseries\small Platform}\\
\endfirsthead
\rowcolor{najdanavy}\textcolor{white}{\sffamily\bfseries\small ID} & \textcolor{white}{\sffamily\bfseries\small Requirement} & \textcolor{white}{\sffamily\bfseries\small Priority} & \textcolor{white}{\sffamily\bfseries\small Platform}\\
\endhead
\bottomrule\endlastfoot
\textbf{FR-PLT-01} & The application is available in English and Arabic with right-to-left layout and light and dark themes; the choice persists. & Must & Web, Mobile\\
\textbf{FR-PLT-02} & A public status page and endpoint report the health of the platform and its dependencies. & Should & Web, Mobile\\
\textbf{FR-PLT-03} & The mobile app checks the latest published release and offers an update. & Should & Mobile\\
\textbf{FR-PLT-04} & The mobile app can be protected by a biometric or device-credential lock. & Could & Mobile\\
\textbf{FR-PLT-05} & A visitor can send a message to the support team from the public support page. & Could & Web\\
\textbf{FR-PLT-06} & Significant administrative and privacy-relevant actions are written to an audit log. & Must & Backend\\
\textbf{FR-PLT-07} & Live changes (incidents, missions, units, chat, vitals, notifications) are pushed to clients over WebSocket so screens update without polling. & Must & Web, Mobile\\
\end{longtable}}


\section{Non-Functional Requirements}\label{sec:nfr}

The non-functional requirements (\cref{tab:nfr}) describe the qualities the system must have. Security,
privacy and reliability dominate, reflecting the nature of the domain.

{\tabfont\rowband
\begin{longtable}{@{}>{\raggedright\arraybackslash}p{\dimexpr(\linewidth-6\tabcolsep)*190/1000\relax}>{\raggedright\arraybackslash}p{\dimexpr(\linewidth-6\tabcolsep)*144/1000\relax}>{\raggedright\arraybackslash}p{\dimexpr(\linewidth-6\tabcolsep)*667/1000\relax}@{}}
\caption{Non-functional requirements}\label{tab:nfr}\\
\rowcolor{najdanavy}\textcolor{white}{\sffamily\bfseries\small ID} & \textcolor{white}{\sffamily\bfseries\small Category} & \textcolor{white}{\sffamily\bfseries\small Requirement}\\
\endfirsthead
\rowcolor{najdanavy}\textcolor{white}{\sffamily\bfseries\small ID} & \textcolor{white}{\sffamily\bfseries\small Category} & \textcolor{white}{\sffamily\bfseries\small Requirement}\\
\endhead
\bottomrule\endlastfoot
\textbf{NFR-SEC-01} & Security & Every API call except a short list of public endpoints is authenticated with a verified Firebase ID token; authorisation is enforced on the server by role (URL rules, method annotations) and by ownership checks, never by client-side hiding.\\
\textbf{NFR-SEC-02} & Security & Secrets (database credentials, Firebase service account, API keys, webhook secrets, field-encryption key, OTP pepper) are supplied by environment configuration and never committed.\\
\textbf{NFR-SEC-03} & Security & Verification codes are stored only as keyed hashes (HMAC-SHA-256); the registration token is one-time and stored hashed; comparisons are constant-time.\\
\textbf{NFR-SEC-04} & Security & Sensitive fields (chat, text evidence, medical free text, emergency contacts, vitals notes) are encrypted at rest at the application layer.\\
\textbf{NFR-SEC-05} & Security & Evidence and documents live in a private bucket; clients receive only short-lived signed upload and download URLs, and file types are restricted by an allow-list.\\
\textbf{NFR-SEC-06} & Security & Abuse-prone public endpoints are rate-limited; the caller-trust policy slows repeat abusers without ever blocking a genuine report.\\
\textbf{NFR-SEC-07} & Security & Third-party webhooks are accepted only with a valid signature inside a freshness window and are processed at most once.\\
\textbf{NFR-PRV-01} & Privacy & A patient’s medical profile is shared only with consent, only with people treating the patient, only while the incident is open, and every access is audited.\\
\textbf{NFR-PRV-02} & Privacy & Users must accept the current Terms and Privacy Policy; versions are immutable and acceptance is recorded with source and time.\\
\textbf{NFR-PRV-03} & Privacy & Live unit positions are visible to staff only; a citizen sees only the units assigned to their incident.\\
\textbf{NFR-PRF-01} & Performance & The submission path does only database work before responding; AI classification, duplicate detection and geocoding run asynchronously after commit.\\
\textbf{NFR-PRF-02} & Performance & Lists that can grow (users, facilities, notifications) are paged and indexed; hot lookups (notifications by user, citizen incidents, active units by facility) have dedicated indexes.\\
\textbf{NFR-PRF-03} & Performance & Live updates travel over a single WebSocket per client rather than repeated polling.\\
\textbf{NFR-REL-01} & Reliability & Retried submissions are safe (idempotent request id, unique index per citizen); the mobile outbox preserves a report across connectivity loss.\\
\textbf{NFR-REL-02} & Reliability & Concurrent actions on the same unit or mission are safe: pessimistic locking on unit assignment, optimistic version checks on incidents and missions.\\
\textbf{NFR-REL-03} & Reliability & Failure of an external service (language model, push, geocoder, mail) degrades only its own feature; core reporting and dispatch continue.\\
\textbf{NFR-REL-04} & Reliability & Cross-system changes that cannot be one transaction (identity provider and database) use compensation so no half-created state remains.\\
\textbf{NFR-USA-01} & Usability & Reporting takes one deliberate action from the home screen; the interface supports English and Arabic with correct right-to-left layout, and light and dark themes.\\
\textbf{NFR-USA-02} & Usability & The system explains refusals in plain language, including the emergency numbers (122 police, 123 ambulance, 180 fire) whenever a report is slowed.\\
\textbf{NFR-MNT-01} & Maintainability & The schema is changed only through versioned Flyway migrations and validated against the entities at start-up.\\
\textbf{NFR-MNT-02} & Maintainability & Business rules that decide safety-critical outcomes are isolated in pure classes with unit tests (caller trust, KYC policy, webhook verification, OTP, facility and role policies).\\
\textbf{NFR-MNT-03} & Maintainability & The backend is organised by feature module with a consistent controller--service--repository layering.\\
\textbf{NFR-OBS-01} & Observability & Health and status endpoints are public; significant events are logged; the audit log records privacy-relevant actions with actor, target and detail.\\
\textbf{NFR-PRT-01} & Portability & Every service ships as a container image or a single archive; configuration is entirely environment-driven, so the stack runs locally with one Compose file.\\
\textbf{NFR-CMP-01} & Compatibility & The web console supports current evergreen browsers; the mobile app supports Android, with the code base cross-platform.\\
\end{longtable}}


\section{User Stories}

User stories express the requirements from the actor's point of view and map to the use cases of the next
section. They are listed by capability area.

\subsection{Identity and access management}{\tabfont\rowband
\begin{longtable}{@{}>{\raggedright\arraybackslash}p{\dimexpr(\linewidth-8\tabcolsep)*83/1000\relax}>{\raggedright\arraybackslash}p{\dimexpr(\linewidth-8\tabcolsep)*143/1000\relax}>{\raggedright\arraybackslash}p{\dimexpr(\linewidth-8\tabcolsep)*417/1000\relax}>{\raggedright\arraybackslash}p{\dimexpr(\linewidth-8\tabcolsep)*357/1000\relax}@{}}
\caption{User stories: identity and access management}\label{tab:us-iam}\\
\rowcolor{najdanavy}\textcolor{white}{\sffamily\bfseries\small ID} & \textcolor{white}{\sffamily\bfseries\small As a\ldots{}} & \textcolor{white}{\sffamily\bfseries\small I want to\ldots{}} & \textcolor{white}{\sffamily\bfseries\small So that\ldots{}}\\
\endfirsthead
\rowcolor{najdanavy}\textcolor{white}{\sffamily\bfseries\small ID} & \textcolor{white}{\sffamily\bfseries\small As a\ldots{}} & \textcolor{white}{\sffamily\bfseries\small I want to\ldots{}} & \textcolor{white}{\sffamily\bfseries\small So that\ldots{}}\\
\endhead
\bottomrule\endlastfoot
\textbf{US-01} & Visitor & register with my e-mail and a code sent to it & my account is tied to an address I really own\\
\textbf{US-02} & Visitor & sign in with my Google account & I do not have to manage another password\\
\textbf{US-03} & New user & be guided to complete my phone, address and gender & responders can reach and find me\\
\textbf{US-04} & User & recover a forgotten password & I am not locked out in an emergency\\
\textbf{US-05} & User & verify or change my phone number & dispatch can call me, and no one else can claim my number\\
\textbf{US-06} & User & change my e-mail address and password & my account stays under my control\\
\textbf{US-07} & Citizen & verify my identity with my ID and a selfie & the service can trust my reports and I can apply as a responder\\
\textbf{US-08} & User & read and accept the current Terms and Privacy Policy & I know how my data is used\\
\textbf{US-09} & Citizen & keep a medical profile and choose whether responders can see it & treatment is safer and I stay in control\\
\textbf{US-10} & Citizen & apply to become a First Responder with my documents & I can help my community\\
\textbf{US-11} & Admin & review applications and their documents & only qualified people become responders\\
\textbf{US-12} & Admin & provision accounts for staff and hand over a one-time link & no password is ever shared\\
\textbf{US-13} & Admin & search users, change roles and facilities, disable accounts & access always matches a person’s job\\
\textbf{US-14} & Super Admin & reconcile Firebase with the database & orphaned or missing accounts are repaired\\
\textbf{US-15} & Super Admin & override identity verification or a phone number when needed & a legitimate user is not stuck\\
\end{longtable}}

\subsection{Emergency reporting}{\tabfont\rowband
\begin{longtable}{@{}>{\raggedright\arraybackslash}p{\dimexpr(\linewidth-8\tabcolsep)*83/1000\relax}>{\raggedright\arraybackslash}p{\dimexpr(\linewidth-8\tabcolsep)*143/1000\relax}>{\raggedright\arraybackslash}p{\dimexpr(\linewidth-8\tabcolsep)*417/1000\relax}>{\raggedright\arraybackslash}p{\dimexpr(\linewidth-8\tabcolsep)*357/1000\relax}@{}}
\caption{User stories: emergency reporting}\label{tab:us-rep}\\
\rowcolor{najdanavy}\textcolor{white}{\sffamily\bfseries\small ID} & \textcolor{white}{\sffamily\bfseries\small As a\ldots{}} & \textcolor{white}{\sffamily\bfseries\small I want to\ldots{}} & \textcolor{white}{\sffamily\bfseries\small So that\ldots{}}\\
\endfirsthead
\rowcolor{najdanavy}\textcolor{white}{\sffamily\bfseries\small ID} & \textcolor{white}{\sffamily\bfseries\small As a\ldots{}} & \textcolor{white}{\sffamily\bfseries\small I want to\ldots{}} & \textcolor{white}{\sffamily\bfseries\small So that\ldots{}}\\
\endhead
\bottomrule\endlastfoot
\textbf{US-16} & Citizen & report an emergency with one tap & help is on its way as fast as possible\\
\textbf{US-17} & Citizen & have my location captured, or place a pin myself & responders find the right spot\\
\textbf{US-18} & Citizen & report on behalf of someone else & I can help a stranger without exposing my own medical data\\
\textbf{US-19} & Citizen & attach photos, video, audio and a message & dispatchers understand the situation\\
\textbf{US-20} & Citizen & add evidence and update the injured count while the incident is live & responders get the latest facts\\
\textbf{US-21} & Citizen & cancel a report I made by mistake & no unit is wasted\\
\textbf{US-22} & Citizen & watch units approach on a map and see their status & I know help is coming\\
\textbf{US-23} & Citizen & chat with the people responding & I can answer questions and give directions\\
\textbf{US-24} & Citizen & review my past incidents & I have a record of what happened\\
\textbf{US-25} & Citizen & read first-aid guidance while I wait & I can act safely until help arrives\\
\textbf{US-26} & Citizen & have my report saved and sent when the connection returns & a poor signal never loses my report\\
\textbf{US-27} & Citizen & be told when to wait and which numbers to call if my reports were slowed & I am never left without a way to get help\\
\end{longtable}}

\subsection{Dispatch operations}{\tabfont\rowband
\begin{longtable}{@{}>{\raggedright\arraybackslash}p{\dimexpr(\linewidth-8\tabcolsep)*83/1000\relax}>{\raggedright\arraybackslash}p{\dimexpr(\linewidth-8\tabcolsep)*143/1000\relax}>{\raggedright\arraybackslash}p{\dimexpr(\linewidth-8\tabcolsep)*417/1000\relax}>{\raggedright\arraybackslash}p{\dimexpr(\linewidth-8\tabcolsep)*357/1000\relax}@{}}
\caption{User stories: dispatch operations}\label{tab:us-dsp}\\
\rowcolor{najdanavy}\textcolor{white}{\sffamily\bfseries\small ID} & \textcolor{white}{\sffamily\bfseries\small As a\ldots{}} & \textcolor{white}{\sffamily\bfseries\small I want to\ldots{}} & \textcolor{white}{\sffamily\bfseries\small So that\ldots{}}\\
\endfirsthead
\rowcolor{najdanavy}\textcolor{white}{\sffamily\bfseries\small ID} & \textcolor{white}{\sffamily\bfseries\small As a\ldots{}} & \textcolor{white}{\sffamily\bfseries\small I want to\ldots{}} & \textcolor{white}{\sffamily\bfseries\small So that\ldots{}}\\
\endhead
\bottomrule\endlastfoot
\textbf{US-28} & Dispatcher & see a live, prioritised queue & the most urgent emergencies are handled first\\
\textbf{US-29} & Dispatcher & see the AI priority and a duplicate suggestion & I triage faster and avoid double dispatch\\
\textbf{US-30} & Dispatcher & see who else is viewing an incident & we do not both assign units\\
\textbf{US-31} & Dispatcher & see available units ranked by distance and ETA & I pick the best unit quickly\\
\textbf{US-32} & Dispatcher & assign several units to one incident & large emergencies get enough help\\
\textbf{US-33} & Dispatcher & cancel an assignment & I can re-route resources\\
\textbf{US-34} & Dispatcher & mark a false report and review caller history & abuse is handled fairly\\
\textbf{US-35} & Dispatcher & chat with the reporter and the crew & everyone shares the same facts\\
\textbf{US-36} & Admin & review flagged callers & repeat abuse is seen and handled by a person\\
\end{longtable}}

\subsection{Field response}{\tabfont\rowband
\begin{longtable}{@{}>{\raggedright\arraybackslash}p{\dimexpr(\linewidth-8\tabcolsep)*83/1000\relax}>{\raggedright\arraybackslash}p{\dimexpr(\linewidth-8\tabcolsep)*143/1000\relax}>{\raggedright\arraybackslash}p{\dimexpr(\linewidth-8\tabcolsep)*417/1000\relax}>{\raggedright\arraybackslash}p{\dimexpr(\linewidth-8\tabcolsep)*357/1000\relax}@{}}
\caption{User stories: field response}\label{tab:us-fld}\\
\rowcolor{najdanavy}\textcolor{white}{\sffamily\bfseries\small ID} & \textcolor{white}{\sffamily\bfseries\small As a\ldots{}} & \textcolor{white}{\sffamily\bfseries\small I want to\ldots{}} & \textcolor{white}{\sffamily\bfseries\small So that\ldots{}}\\
\endfirsthead
\rowcolor{najdanavy}\textcolor{white}{\sffamily\bfseries\small ID} & \textcolor{white}{\sffamily\bfseries\small As a\ldots{}} & \textcolor{white}{\sffamily\bfseries\small I want to\ldots{}} & \textcolor{white}{\sffamily\bfseries\small So that\ldots{}}\\
\endhead
\bottomrule\endlastfoot
\textbf{US-37} & Crew member & start or join a shift on my unit & dispatch knows we are available\\
\textbf{US-38} & Unit LEAD & be alerted loudly when a mission is offered, and accept or reject it & assignments are acknowledged immediately\\
\textbf{US-39} & Unit LEAD & mark en route, arrived and complete & everyone sees progress\\
\textbf{US-40} & Responder & see a route to the incident and my teammates & we arrive together\\
\textbf{US-41} & Unit LEAD & share my unit’s live position & dispatch and the citizen can track us\\
\textbf{US-42} & Responder & review my past missions & I can account for my work\\
\textbf{US-43} & Admin & force-end a stuck shift & a unit can be reassigned\\
\end{longtable}}

\subsection{Hospital coordination}{\tabfont\rowband
\begin{longtable}{@{}>{\raggedright\arraybackslash}p{\dimexpr(\linewidth-8\tabcolsep)*83/1000\relax}>{\raggedright\arraybackslash}p{\dimexpr(\linewidth-8\tabcolsep)*143/1000\relax}>{\raggedright\arraybackslash}p{\dimexpr(\linewidth-8\tabcolsep)*417/1000\relax}>{\raggedright\arraybackslash}p{\dimexpr(\linewidth-8\tabcolsep)*357/1000\relax}@{}}
\caption{User stories: hospital coordination}\label{tab:us-hos}\\
\rowcolor{najdanavy}\textcolor{white}{\sffamily\bfseries\small ID} & \textcolor{white}{\sffamily\bfseries\small As a\ldots{}} & \textcolor{white}{\sffamily\bfseries\small I want to\ldots{}} & \textcolor{white}{\sffamily\bfseries\small So that\ldots{}}\\
\endfirsthead
\rowcolor{najdanavy}\textcolor{white}{\sffamily\bfseries\small ID} & \textcolor{white}{\sffamily\bfseries\small As a\ldots{}} & \textcolor{white}{\sffamily\bfseries\small I want to\ldots{}} & \textcolor{white}{\sffamily\bfseries\small So that\ldots{}}\\
\endhead
\bottomrule\endlastfoot
\textbf{US-44} & Ambulance LEAD & choose the nearest suitable hospital & the patient is received without delay\\
\textbf{US-45} & Ambulance LEAD & send the patient’s vital signs during transport & the hospital prepares in advance\\
\textbf{US-46} & Hospital staff & see inbound patients, ETA and live vitals & we get the right team ready\\
\textbf{US-47} & Treating staff & see the patient’s medical profile when they consented & we avoid dangerous mistakes\\
\end{longtable}}

\subsection{Administration and platform}{\tabfont\rowband
\begin{longtable}{@{}>{\raggedright\arraybackslash}p{\dimexpr(\linewidth-8\tabcolsep)*83/1000\relax}>{\raggedright\arraybackslash}p{\dimexpr(\linewidth-8\tabcolsep)*143/1000\relax}>{\raggedright\arraybackslash}p{\dimexpr(\linewidth-8\tabcolsep)*417/1000\relax}>{\raggedright\arraybackslash}p{\dimexpr(\linewidth-8\tabcolsep)*357/1000\relax}@{}}
\caption{User stories: administration and platform}\label{tab:us-adm}\\
\rowcolor{najdanavy}\textcolor{white}{\sffamily\bfseries\small ID} & \textcolor{white}{\sffamily\bfseries\small As a\ldots{}} & \textcolor{white}{\sffamily\bfseries\small I want to\ldots{}} & \textcolor{white}{\sffamily\bfseries\small So that\ldots{}}\\
\endfirsthead
\rowcolor{najdanavy}\textcolor{white}{\sffamily\bfseries\small ID} & \textcolor{white}{\sffamily\bfseries\small As a\ldots{}} & \textcolor{white}{\sffamily\bfseries\small I want to\ldots{}} & \textcolor{white}{\sffamily\bfseries\small So that\ldots{}}\\
\endhead
\bottomrule\endlastfoot
\textbf{US-48} & Admin & create, edit and seed facilities, and see the impact of a type change & the map of facilities is accurate\\
\textbf{US-49} & Admin & create, edit and retire response units and reconcile their status & the fleet in the system matches the street\\
\textbf{US-50} & Facility Manager & manage my facility’s units and staff & I run my station without needing an administrator\\
\textbf{US-51} & Admin & edit first-aid guidance & advice stays current\\
\textbf{US-52} & Super Admin & publish a new version of the Terms or Privacy Policy and require re-acceptance & legal changes take effect properly\\
\textbf{US-53} & User & receive notifications in my language, in the app and as push & I never miss something that matters\\
\textbf{US-54} & User & switch language and theme & the app is comfortable for me\\
\textbf{US-55} & Mobile user & be offered app updates and lock the app with biometrics & I have the latest fixes and my data stays private\\
\textbf{US-56} & Visitor & see the system status and contact support & I know if the service is healthy and can get help\\
\end{longtable}}


\section{Use Case Analysis}\label{sec:usecases}

The system's behaviour is specified by 64 use cases in seven groups. This section presents the use-case diagram of
each group; the complete specification of every use case --- actor, platform, preconditions, postconditions, main
flow and alternate flows --- is in Appendix~\ref{app:use-cases}. Three of the most important use cases are summarised
afterwards. In the diagrams, blue-shaded use cases are restricted to administrators, and actors marked
\guillemotleft system\guillemotright{} are external systems.

\subsection{Account, access and identity}
\Cref{fig:uc-iam-a} covers self-service: registration by e-mail code or Google, profile completion, sign-in,
account changes, legal consent, identity verification and the medical profile. \Cref{fig:uc-iam-b} covers the
administrative side: provisioning staff, administering users, reviewing responder applications, and the
Super Admin's reconciliation and override tools.
\diagram{uc_iam_a.png}{Use cases for registration, sign-in, profile, identity verification and medical profile.}{fig:uc-iam-a}
\diagram{uc_iam_b.png}{Use cases for account administration and responder applications.}{fig:uc-iam-b}

\subsection{Emergency reporting}
Reporting is the heart of the citizen's experience (\cref{fig:uc-report}). Submitting a report always includes a
location check and a consent-and-identity gate; sending on behalf of someone else, sharing the medical profile,
attaching evidence and queueing offline are optional extensions.
\diagram{uc_report.png}{Use cases for emergency reporting.}{fig:uc-report}

\subsection{Dispatch operations}
Dispatchers (and the administrators, who inherit their capabilities) monitor the live queue, review the AI
priority, resolve duplicates, assign and cancel units, and handle false reports (\cref{fig:uc-dispatch}).
\diagram{uc_dispatch.png}{Use cases for dispatch operations.}{fig:uc-dispatch}

\subsection{Field response and shift management}
Field staff work on response units: a shift ties employees to a unit with one \emph{lead}, who is accountable
for the mission, and any number of \emph{crew} (\cref{fig:uc-field}).
\diagram{uc_field.png}{Use cases for field response and shift management.}{fig:uc-field}

\subsection{Hospital coordination}
An ambulance lead chooses a receiving hospital, records vital signs during transport and completes the mission by
handing the patient over; the hospital monitors incoming transfers and, where the patient has consented, sees the
medical profile (\cref{fig:uc-hospital}).
\diagram{uc_hospital.png}{Use cases for hospital coordination.}{fig:uc-hospital}

\subsection{Facility, unit and content administration}
Administrators manage facilities and units and the content shown to users; facility managers run the units and staff
of their own facility; the Super Admin publishes the legal documents (\cref{fig:uc-admin}).
\diagram{uc_admin.png}{Use cases for facility, unit and content administration.}{fig:uc-admin}

\subsection{Notifications and general services}
\Cref{fig:uc-general} covers notifications, push-device registration, language and theme, support, system status,
the app-update check and the biometric lock.
\diagram{uc_general.png}{Use cases for notifications and general services.}{fig:uc-general}

\subsection{Key use cases in brief}

\paragraph{UC-20: Report an emergency.} The citizen chooses a category; the app captures the location and
submits the report with a client-generated request id. The server replays an existing incident for the same id,
otherwise validates the location, checks that the citizen has no other active incident and is not cooling down,
decides medical-profile sharing, saves the incident and starts AI priority, duplicate detection and reverse
geocoding asynchronously. After commit, dispatchers receive a live update and an urgent notification.

\paragraph{UC-32 and UC-39: Assign a unit and work the mission.} The dispatcher selects an available unit. The
server locks the unit, checks that it is available with an active lead and that the incident is open, and creates
a mission in the \emph{offered} state. The unit's crew receive an urgent notification; the lead accepts, marks the
unit en route and arrived, and completes the mission. Completing the last active mission resolves the incident.

\paragraph{UC-43 to UC-45: Hospital handoff.} When an ambulance has arrived on scene, its lead selects a registered
hospital within 30\,km. The hospital's staff are alerted; the lead marks the transfer en route, records vital signs
during transport, and marks it arrived, which completes the mission, frees the unit and resolves the incident if
it was the last active mission.

\section{Overall System Workflow}

\Cref{fig:workflow} shows the whole life of an emergency across the five parties as an activity diagram with
swim lanes. \Cref{fig:validation} zooms in on the ordered checks that every report passes before it is accepted;
the order matters --- consent and identity first, then idempotent replay, then location, then the active-incident
rule and finally the caller-trust cool-down.

\diagram{act_lifecycle.png}{End-to-end handling of an emergency, from report to resolution.}{fig:workflow}
\diagramwide{act_validate.png}{The checks a report passes, in order, before it is accepted.}{fig:validation}

\section{Business Rules}\label{sec:business-rules}

Business rules are the constraints that keep the service safe and fair. They are enforced on the server, in the
layer named in the last column of each table.

\subsection{Accounts, roles and verification}{\tabfont\rowband
\begin{longtable}{@{}>{\raggedright\arraybackslash}p{\dimexpr(\linewidth-6\tabcolsep)*92/1000\relax}>{\raggedright\arraybackslash}p{\dimexpr(\linewidth-6\tabcolsep)*736/1000\relax}>{\raggedright\arraybackslash}p{\dimexpr(\linewidth-6\tabcolsep)*172/1000\relax}@{}}
\caption{Business rules: accounts, roles and verification}\label{tab:br-accounts}\\
\rowcolor{najdanavy}\textcolor{white}{\sffamily\bfseries\small ID} & \textcolor{white}{\sffamily\bfseries\small Rule} & \textcolor{white}{\sffamily\bfseries\small Enforced in}\\
\endfirsthead
\rowcolor{najdanavy}\textcolor{white}{\sffamily\bfseries\small ID} & \textcolor{white}{\sffamily\bfseries\small Rule} & \textcolor{white}{\sffamily\bfseries\small Enforced in}\\
\endhead
\bottomrule\endlastfoot
\textbf{BR-1} & A user has exactly one role. E-mail addresses are unique; a phone number can be verified on at most one account. & Database, user service\\
\textbf{BR-2} & E-mail registration: six-digit code valid 10 minutes; 60 seconds between codes; five wrong attempts per code; five codes per e-mail per hour; ten code requests, ten registrations and thirty verifications (per 10 minutes) per client address. A verified e-mail stays usable for 15 minutes through a one-time token. & Auth controller, OTP service\\
\textbf{BR-3} & Registration is atomic across Firebase and the database; on database failure the Firebase account is deleted. & Auth service\\
\textbf{BR-4} & A profile is complete when phone, governorate, city, street and gender are present. Only the server computes this. & Profile evaluator\\
\textbf{BR-5} & Roles that must be linked to a facility: hospital staff (hospital), ambulance crew (ambulance station or hospital), police (police station), firefighter (fire station), facility manager (any type). Other roles cannot be linked to a facility. & Role--facility policy\\
\textbf{BR-6} & Only a Super Admin may create, grant or modify a Super Admin; a Super Admin cannot remove their own role; the only remaining Super Admin cannot be disabled; a plain Admin cannot disable a Super Admin. & Role and user services\\
\textbf{BR-7} & A disabled account is refused on every request. Re-enabling clears the deletion mark. & Token filter, user service\\
\textbf{BR-8} & Employees get a temporary password they never see; they set their own through a one-time reset link. & Auth service\\
\textbf{BR-9} & A citizen may apply to become a First Responder only with verified e-mail and phone, a VERIFIED identity and no pending application, with a motivation and files. Approval needs at least one document; it changes the role, updates the claim and creates the personal unit (one per responder). & Application service\\
\textbf{BR-10} & Identity-verification status is written only by the backend (signed webhook confirmed with the provider, provider reconciliation, or an audited Super Admin action). A settled status is not undone by a late “in progress” message. & KYC service\\
\textbf{BR-11} & When required by configuration, a citizen must be VERIFIED to report. Admin roles are exempt by default and accounts created before a configured date are grandfathered. & KYC policy, incident controller\\
\textbf{BR-12} & Legal versions are immutable and numbered per document type. A new version that requires acceptance makes every protected call return 403 LEGAL\_\allowbreak{}CONSENT\_\allowbreak{}REQUIRED until accepted. & Legal services, interceptor\\
\end{longtable}}

\subsection{Reporting}{\tabfont\rowband
\begin{longtable}{@{}>{\raggedright\arraybackslash}p{\dimexpr(\linewidth-6\tabcolsep)*92/1000\relax}>{\raggedright\arraybackslash}p{\dimexpr(\linewidth-6\tabcolsep)*736/1000\relax}>{\raggedright\arraybackslash}p{\dimexpr(\linewidth-6\tabcolsep)*172/1000\relax}@{}}
\caption{Business rules: reporting}\label{tab:br-reporting}\\
\rowcolor{najdanavy}\textcolor{white}{\sffamily\bfseries\small ID} & \textcolor{white}{\sffamily\bfseries\small Rule} & \textcolor{white}{\sffamily\bfseries\small Enforced in}\\
\endfirsthead
\rowcolor{najdanavy}\textcolor{white}{\sffamily\bfseries\small ID} & \textcolor{white}{\sffamily\bfseries\small Rule} & \textcolor{white}{\sffamily\bfseries\small Enforced in}\\
\endhead
\bottomrule\endlastfoot
\textbf{BR-13} & A citizen may have only one active incident (NEW, AI\_\allowbreak{}PROCESSED, DISPATCHER\_\allowbreak{}REVIEW, ASSIGNED). & Incident service\\
\textbf{BR-14} & A repeated submission with the same client request id returns the original incident before any other rule is applied and causes no side effects; the database enforces uniqueness per citizen. & Incident service, unique index\\
\textbf{BR-15} & The incident location must be inside Egypt. A manual pin is allowed only when switched on and within the configured radius (default 300 m) of the device position, which must also be in Egypt. & Location policy\\
\textbf{BR-16} & Someone else’s emergency never carries the reporter’s medical profile. For self-reports the profile is attached only if the citizen allows it for this report and the standing consent is on. & Incident service\\
\textbf{BR-17} & The citizen may update the injured count while the incident is active. Until a unit is assigned the citizen may also edit the original message and delete evidence; evidence can be added while the incident is live. & Incident and media services\\
\textbf{BR-18} & Cancelling needs a reason category; “other” needs an explanation (length-limited). It is refused once a unit has arrived. It stands down every offered, accepted or en-route mission and frees the units. & Incident service\\
\textbf{BR-19} & A false-report verdict (prank or hoax, spam or test, malicious) can be set once by dispatch staff and cleared by an Admin. The incident is closed and units stood down only if no unit has arrived and it is still open; otherwise the verdict is recorded and the status left alone. Both actions are audited. & Incident service\\
\textbf{BR-20} & A duplicate suggestion considers open incidents of the same category within 2 km and 120 minutes. Only a dispatcher confirms: the duplicate closes as cancelled, linked to the original. & AI duplicate service, incident service\\
\textbf{BR-21} & AI failures never affect an incident. AI is disabled by a switch and automatically when no key is configured. & AI services\\
\end{longtable}}

\subsection{Dispatch and field response}{\tabfont\rowband
\begin{longtable}{@{}>{\raggedright\arraybackslash}p{\dimexpr(\linewidth-6\tabcolsep)*92/1000\relax}>{\raggedright\arraybackslash}p{\dimexpr(\linewidth-6\tabcolsep)*736/1000\relax}>{\raggedright\arraybackslash}p{\dimexpr(\linewidth-6\tabcolsep)*172/1000\relax}@{}}
\caption{Business rules: dispatch and field response}\label{tab:br-dispatch}\\
\rowcolor{najdanavy}\textcolor{white}{\sffamily\bfseries\small ID} & \textcolor{white}{\sffamily\bfseries\small Rule} & \textcolor{white}{\sffamily\bfseries\small Enforced in}\\
\endfirsthead
\rowcolor{najdanavy}\textcolor{white}{\sffamily\bfseries\small ID} & \textcolor{white}{\sffamily\bfseries\small Rule} & \textcolor{white}{\sffamily\bfseries\small Enforced in}\\
\endhead
\bottomrule\endlastfoot
\textbf{BR-22} & A unit can be offered a mission only if it is AVAILABLE, not archived, has an active LEAD on shift, and the incident is not closed. Assignment takes a pessimistic lock on the unit so two dispatchers cannot book it together. & Mission service\\
\textbf{BR-23} & An incident can have several missions. The incident becomes ASSIGNED with the first. & Mission service\\
\textbf{BR-24} & A unit whose current crew includes the person who reported the emergency cannot be assigned to it --- the reporter is the one who needs help. When a mission is offered, everyone on the unit’s active shift is recorded as a participant. & Mission service\\
\textbf{BR-25} & Only the unit’s LEAD accepts, rejects, advances (en route, arrived), completes or withdraws a mission. Withdrawing is allowed only from accepted or en route. Dispatch can cancel at any non-terminal stage. & Mission service\\
\textbf{BR-26} & Rejected, cancelled and completed missions free the unit. When the last active mission is cancelled or withdrawn while the incident is ASSIGNED, the incident returns to DISPATCHER\_\allowbreak{}REVIEW; when the last active mission completes, the incident is RESOLVED. & Mission service\\
\textbf{BR-27} & A unit has at most one active LEAD; an employee is on one shift at a time; nobody can end or leave a shift while the unit is on an active mission. Starting the LEAD shift sets the unit AVAILABLE; ending it sets it OFFLINE; a crew member leaving changes nothing. Archived units cannot start shifts. & Shift service\\
\textbf{BR-28} & Only the LEAD updates a unit’s position; positions are visible to staff only; the LEAD may reset the position to the home facility. & Unit controller, service\\
\textbf{BR-29} & Chat: the citizen, dispatch staff and participants of an active mission (offered to arrived) may send; anyone who ever participated may read. Chat notifications never include the message text. & Chat service, notification publisher\\
\textbf{BR-30} & Concurrent edits are guarded by optimistic version fields on incidents and missions. & JPA\\
\end{longtable}}


\subsection{Caller trust}\label{sec:caller-trust}

The caller-trust policy exists to slow repeated abuse without ever refusing a genuine emergency. It is a pure,
deterministic function of a caller's recent history. Each false-report verdict and each late cancellation adds a
weight to a score; the weight fades with a half-life of 30 days and the event stops counting after 120 days:
\begin{equation}
  s \;=\; \sum_{e} w_e \cdot 2^{-\Delta t_e/30}, \qquad \Delta t_e \le 120\ \text{days},
\end{equation}
where $\Delta t_e$ is the age of event $e$ in days. The score determines a tier, and the tier a minimum gap between
reports.

\begin{table}[!htbp]
\caption{Parameters of the caller-trust policy.}\label{tab:trust-params}
\centering\tabfont\rowband
\begin{tabular}{@{}>{\bfseries}p{3.4cm}p{10.6cm}@{}}
\rowcolor{najdanavy}\textcolor{white}{\sffamily\bfseries\small Item} & \textcolor{white}{\sffamily\bfseries\small Value} \\
False-report weights & Malicious 10; prank or hoax 6; spam or test 3. \\
Cancellation weights & Zero unless a unit had already accepted; then submitted by mistake or other 2; situation resolved, got help elsewhere or duplicate 0.5. \\
Decay & Each event halves every 30 days and stops counting after 120 days. \\
Tiers & \code{NONE} below 3; \code{WATCH} from 3 (appears on the admin review list); \code{COOLDOWN} from 6; \code{REVIEW} from 10. \\
Effect & \code{COOLDOWN} and \code{REVIEW} require 10 minutes between new reports. A report inside the wait receives HTTP 429 with a \code{Retry-After} header and a message naming 122, 123 and 180. Nothing else is ever refused automatically; suspension is a manual administrator decision. \\
\end{tabular}
\end{table}

\begin{keyidea}[An emergency line must fail open]
No automatic rule in \ProjectName{} ever refuses a report. The harshest effect is a ten-minute gap between reports,
together with the emergency numbers, and a human decides whether to suspend an account. Abuse is handled by delay
and by human verdicts, not by denial.
\end{keyidea}

\subsection{Hospital coordination}{\tabfont\rowband
\begin{longtable}{@{}>{\raggedright\arraybackslash}p{\dimexpr(\linewidth-6\tabcolsep)*92/1000\relax}>{\raggedright\arraybackslash}p{\dimexpr(\linewidth-6\tabcolsep)*736/1000\relax}>{\raggedright\arraybackslash}p{\dimexpr(\linewidth-6\tabcolsep)*172/1000\relax}@{}}
\caption{Business rules: hospital coordination}\label{tab:br-hospital}\\
\rowcolor{najdanavy}\textcolor{white}{\sffamily\bfseries\small ID} & \textcolor{white}{\sffamily\bfseries\small Rule} & \textcolor{white}{\sffamily\bfseries\small Enforced in}\\
\endfirsthead
\rowcolor{najdanavy}\textcolor{white}{\sffamily\bfseries\small ID} & \textcolor{white}{\sffamily\bfseries\small Rule} & \textcolor{white}{\sffamily\bfseries\small Enforced in}\\
\endhead
\bottomrule\endlastfoot
\textbf{BR-31} & Only the LEAD of an ambulance mission that is ARRIVED can select a destination. Exactly one of a registered hospital or a free-text destination is required; one transfer per mission. & Hospital transfer service\\
\textbf{BR-32} & The recommended list contains registered hospitals within 30 km of the incident, nearest first. & Hospital transfer service\\
\textbf{BR-33} & Selecting a registered hospital notifies its staff. A free-text destination notifies nobody and accepts no vitals. & Transfer service, notification publisher\\
\textbf{BR-34} & Transfer states advance SELECTED $\rightarrow$ EN\_\allowbreak{}ROUTE $\rightarrow$ ARRIVED in order; arrival completes the mission, frees the unit and resolves the incident if it was the last active mission. & Transfer service\\
\textbf{BR-35} & A medical profile is visible to a dispatcher; to the crew of an ambulance or first-responder unit on an active mission; and to staff of the hospital selected as the destination, from selection until the mission closes. Only while the incident is open, only if consent is on, never to administrators; every read is audited and fails if the audit cannot be written. & Incident medical-profile service\\
\end{longtable}}

\subsection{Facilities, units and content}{\tabfont\rowband
\begin{longtable}{@{}>{\raggedright\arraybackslash}p{\dimexpr(\linewidth-6\tabcolsep)*92/1000\relax}>{\raggedright\arraybackslash}p{\dimexpr(\linewidth-6\tabcolsep)*736/1000\relax}>{\raggedright\arraybackslash}p{\dimexpr(\linewidth-6\tabcolsep)*172/1000\relax}@{}}
\caption{Business rules: facilities, units and content}\label{tab:br-facilities}\\
\rowcolor{najdanavy}\textcolor{white}{\sffamily\bfseries\small ID} & \textcolor{white}{\sffamily\bfseries\small Rule} & \textcolor{white}{\sffamily\bfseries\small Enforced in}\\
\endfirsthead
\rowcolor{najdanavy}\textcolor{white}{\sffamily\bfseries\small ID} & \textcolor{white}{\sffamily\bfseries\small Rule} & \textcolor{white}{\sffamily\bfseries\small Enforced in}\\
\endhead
\bottomrule\endlastfoot
\textbf{BR-36} & A facility is “registered” while it has an active linked staff member or a non-archived unit. The flag is refreshed on every link change and reconciled nightly. Only registered hospitals can be destinations. & Facility service, reconciler\\
\textbf{BR-37} & A facility manager runs exactly one facility; the unit type (hospital or ambulance station $\rightarrow$ ambulance, fire station $\rightarrow$ fire truck, police station $\rightarrow$ police car) and the permitted staff roles come from the facility type, never from the request. & Facility-manager policy\\
\textbf{BR-38} & Plate numbers are unique. A unit with no history is deleted; a unit with history is archived (offline, kept for the record). Removal is refused while the unit is on a live mission or someone is on shift with it. & Facility-manager service\\
\textbf{BR-39} & Changing a facility’s type shows an impact assessment first; units and staff that no longer fit are handled by the change. & Facility type-change service\\
\textbf{BR-40} & Facility import skips unnamed entries and facilities already at the same coordinates; ambulance stations are never imported. & Facility seeder\\
\end{longtable}}

\subsection{Data protection and retention}{\tabfont\rowband
\begin{longtable}{@{}>{\raggedright\arraybackslash}p{\dimexpr(\linewidth-6\tabcolsep)*92/1000\relax}>{\raggedright\arraybackslash}p{\dimexpr(\linewidth-6\tabcolsep)*736/1000\relax}>{\raggedright\arraybackslash}p{\dimexpr(\linewidth-6\tabcolsep)*172/1000\relax}@{}}
\caption{Business rules: data protection and retention}\label{tab:br-data}\\
\rowcolor{najdanavy}\textcolor{white}{\sffamily\bfseries\small ID} & \textcolor{white}{\sffamily\bfseries\small Rule} & \textcolor{white}{\sffamily\bfseries\small Enforced in}\\
\endfirsthead
\rowcolor{najdanavy}\textcolor{white}{\sffamily\bfseries\small ID} & \textcolor{white}{\sffamily\bfseries\small Rule} & \textcolor{white}{\sffamily\bfseries\small Enforced in}\\
\endhead
\bottomrule\endlastfoot
\textbf{BR-41} & Chat content, text evidence, free-text medical fields, emergency-contact names and phones, and vitals notes are encrypted at rest with a configured key. & Attribute converters\\
\textbf{BR-42} & Evidence and application documents live in a private bucket; access is by signed URL only; incident deletion removes the evidence folder after commit. & Storage service\\
\textbf{BR-43} & Notifications and device tokens older than 90 days are purged nightly; webhook event ids after 7 days; expired verification rows every 15 minutes. & Scheduled jobs\\
\textbf{BR-44} & Medical-profile reads, false-report verdicts and their clearing, phone and e-mail overrides, employee creation, unit deletion or archiving and identity overrides are audited. & Services\\
\end{longtable}}


\section{Assumptions and Constraints}

The system is built on assumptions about its environment (for example, that users have a smartphone with a
location sensor, and that dispatch desks have continuous access to the console) and operates within constraints
(for example, that reports are accepted only for locations inside Egypt and that the mobile app is built only for
Android). Both are listed in Appendix~\ref{app:decisions}.

\section{Requirements Traceability}\label{sec:traceability}

\Cref{tab:rtm} links each group of requirements to the use cases that realise them, the modules that implement
them, and the automated tests and scenarios that verify them (\cref{chap:testing}).

{\tabfont\rowband
\begin{longtable}{@{}>{\raggedright\arraybackslash}p{\dimexpr(\linewidth-8\tabcolsep)*211/1000\relax}>{\raggedright\arraybackslash}p{\dimexpr(\linewidth-8\tabcolsep)*158/1000\relax}>{\raggedright\arraybackslash}p{\dimexpr(\linewidth-8\tabcolsep)*316/1000\relax}>{\raggedright\arraybackslash}p{\dimexpr(\linewidth-8\tabcolsep)*316/1000\relax}@{}}
\caption{Requirements traceability matrix}\label{tab:rtm}\\
\rowcolor{najdanavy}\textcolor{white}{\sffamily\bfseries\small Requirements} & \textcolor{white}{\sffamily\bfseries\small Use cases} & \textcolor{white}{\sffamily\bfseries\small Implemented in} & \textcolor{white}{\sffamily\bfseries\small Verified by}\\
\endfirsthead
\rowcolor{najdanavy}\textcolor{white}{\sffamily\bfseries\small Requirements} & \textcolor{white}{\sffamily\bfseries\small Use cases} & \textcolor{white}{\sffamily\bfseries\small Implemented in} & \textcolor{white}{\sffamily\bfseries\small Verified by}\\
\endhead
\bottomrule\endlastfoot
\textbf{FR-IAM-01 -- 04} & UC-1 & auth, mail, security & TC-01 -- 03\\
\textbf{FR-IAM-05 -- 11} & UC-2 -- 8 & auth, user & TC-01, TC-28\\
\textbf{FR-IAM-12 -- 18} & UC-12 -- 15, 18 & auth, user, security & TC-04, TC-05\\
\textbf{FR-IAM-19 -- 20} & UC-16, 17 & user, notification & Application scenario\\
\textbf{FR-ID-01 -- 05} & UC-10, 19 & kyc & KycPolicyTest, KycStatusTest, DiditWebhookVerifierTest, TC-06\\
\textbf{FR-LEG-01 -- 03} & UC-9, 57 & legal & TC-07\\
\textbf{FR-MED-01 -- 05} & UC-11, 47 & user, incident & medical\_\allowbreak{}profile\_\allowbreak{}test, TC-12 -- 14\\
\textbf{FR-REP-01 -- 04, 14} & UC-20 & incident, geo & geo\_\allowbreak{}boundary\_\allowbreak{}test, TC-08 -- 10\\
\textbf{FR-REP-05} & UC-20 & incident & TC-12\\
\textbf{FR-REP-06 -- 07} & UC-21, 22 & incident (media, storage) & Evidence scenario\\
\textbf{FR-REP-08} & UC-23 & incident & TC-19\\
\textbf{FR-REP-09 -- 11} & UC-24 -- 26 & incident, WebSocket publishers & Tracking scenario\\
\textbf{FR-REP-12} & UC-27, 56 & guidance, content & Guidance scenario\\
\textbf{FR-REP-13} & UC-28 & mobile outbox & outbox\_\allowbreak{}report\_\allowbreak{}test, TC-11\\
\textbf{FR-DSP-01 -- 05} & UC-29 -- 32 & incident, ai, web console & TC-29\\
\textbf{FR-DSP-06 -- 09} & UC-32 -- 34 & incident (missions) & TC-15 -- 18\\
\textbf{FR-DSP-10 -- 13} & UC-35, 36, 55 & incident (false report, caller trust) & CallerTrustPolicyTest, TC-20, 21\\
\textbf{FR-DSP-14} & UC-37 & incident & Deletion scenario\\
\textbf{FR-FLD-01 -- 04} & UC-38, 42 & unit (shifts) & Shift scenario\\
\textbf{FR-FLD-05 -- 10} & UC-39 -- 41 & incident (missions), unit & TC-17\\
\textbf{FR-HOS-01 -- 06} & UC-43 -- 46 & incident (transfers), hospital & TC-22, 23\\
\textbf{FR-ADM-01 -- 04} & UC-48 -- 50 & facility & FacilityCsvParserTest, FacilityCsvSeederTest\\
\textbf{FR-ADM-05 -- 06} & UC-51, 52 & unit & TC-25\\
\textbf{FR-ADM-07} & UC-55, 56 & incident, guidance & ---\\
\textbf{FR-ADM-08} & UC-53, 54 & unit (facility manager) & FacilityManagerPolicyTest, TC-24\\
\textbf{FR-NOT-01 -- 05} & UC-58, 59 & notification & NotificationMessagesTest, TC-26, 27\\
\textbf{FR-PLT-01} & UC-60 & clients & TC-28\\
\textbf{FR-PLT-02 -- 05} & UC-61 -- 64 & status, appversion, clients & AppReleaseServiceTest, ReleaseVersionsTest, TC-30\\
\textbf{FR-PLT-06 -- 07} & --- & audit, WebSocket & Audit scenarios in TC-14, TC-20\\
\end{longtable}}


\section{Chapter Summary}

This chapter defined \ProjectName{} through its ten roles, 97 functional requirements in nine capability groups,
25 non-functional requirements, 56 user stories, 64 use cases and the business rules that keep the service safe.
The qualities that distinguish it --- fail-open abuse control, consent-based medical data, idempotent reporting and
graceful degradation --- are stated as requirements here and realised in \cref{chap:design,chap:implementation}.
